% Folder 147, batch 06 (archivists' pages 101-120).
% Pass: Opus 5.5 (claude-opus-5-5), 2026-10-09 — first pass, unchecked against the pages by a human.
%
\documentclass[11pt,a4paper]{article}
\input{../preamble/grothendieck}

\folder{147}
\batch{6}
\pages{101}{120}
\foldertitle{Structure à l'infini des Mg,ν [pages 1 à 90, dont table des matières] : notes manuscrites (s.d.).}
\dating{s.d. — le groupe « Autour de La "Longue Marche" à travers la théorie de Galois » (141 à 149) est daté [à partir de 1978]-1983}
\watermark{Édition de démonstration}

\begin{document}

\page{101}
\note{p. 50 de l'auteur. La phrase commence à la page précédente, hors du lot.}
morphismes de normalisation --- il transforme morph.\ de norm.\ stricts en
\struck{isomorphismes} \add{équivalences} de \struck{\ill{}} groupoïdes, et
\uncertain{sommes} de \ill{} stables en produit des groupoïdes correspondants.

Pour nous, l'intérêt principal est dans la formule (obtenue en prenant le
« sous-objet strict » minimal de $G$)
\[
  \text{(144)}\qquad \Pi_1 M_G \;\simeq\; \prod_{\alpha\in S} \Pi_1 M_{g_\alpha, \hat{I}_\alpha}
\]
\note{l'indice $g_\alpha$ est récrit sur une autre lettre, et un exposant biffé suit la parenthèse fermante.}
qui tient compte \add{d'}\uncertain{ailleurs} (et redonne \uncertain{tautologiquement}) les
\struck{isomorphismes} \add{équivalences} des $\Pi M_{G'} \to \Pi M_G$ par des
morphismes de norm.\ stricts (qui induisent en effet des \uncertain{isom.} pour les
normalisés stricts $\widetilde{G}'$ \uncertain{et} $\widetilde{G}$ \uncertain{de} $G$,

\begin{tikzcd}
  \widetilde{G}' \arrow[r] \arrow[d] & \widetilde{G} \arrow[d] \\
  G' \arrow[r] & G
\end{tikzcd}

\note{ce carré porte le numéro (145).}
$G' \to G$ est un morph.\ de norm.\ stricts ssi le morphisme correspondant sur
les normalisés (\uncertain{totalement}) est un \emph{iso} …)

\page{102}
Mais il faut encore expliciter la structure de $S\Pi M_G$, en termes des
$M_{g_\alpha, \hat{I}_\alpha}$ --- nous le ferons au n° suivant.

Pour un morphisme de normalisation quelconque (pas nécessairement strict), le
diagramme commutatif donne un diagramme commun (à \ill{} près)

\begin{tikzcd}[column sep=small, row sep=small, nodes={font=\scriptsize}]
  \Pi_1 M_G \arrow[rr] \arrow[d, "\wr"] & & \Pi_1 M_{G'} \arrow[d, "\wr"] \\
  {\Pi_1 M_{\widetilde{G}} = \prod_{\alpha\in S} M_{g_\alpha, \hat{I}_\alpha}} \arrow[dr] & & {\Pi_1 M_{\widetilde{G}'} = \prod_{\alpha'\in S'} \Pi_1 M_{g'_{\alpha'}, \hat{I}'_{\alpha'}}} \\
  & {\prod_{\alpha'\in S'} M_{g_\alpha, \hat{I}_\alpha}} \arrow[ur, "\Psi"'] &
\end{tikzcd}

\note{ce diagramme porte le numéro (146) ; l'indice de $M$ en haut à gauche est récrit sur un autre. La flèche oblique descendante porte l'annotation : « en identifiant $S'$ à une partie de $S$, d'où $g_\alpha = g'_\alpha$ » (lecture incertaine).}
\marginal{Ceci provient d'un diagramme analogue au niveau des multiplicités modulaires …\ (lecture incertaine, écrit en oblique dans la marge gauche)}

La flèche oblique $\Psi$ est induite par des flèches
\[
  \text{(147)}\qquad \Pi_1 M_{g_\alpha, \hat{I}_\alpha} \longrightarrow \Pi_1 M_{g'_\alpha, \hat{I}'_\alpha}
\]
compte tenu du fait qu'on a $\hat{I}'_\alpha \subset \hat{I}_\alpha$.
Ainsi la connaissance des morphismes de normalisation \add{générale} \ill{} se
réduit à celle des morphismes
\[
  \text{(148)}\qquad \Pi_1 M_{g,I} \longrightarrow \Pi_1 M_{g,I'} \qquad \text{pour } I' \subset I,
\]

\page{103}
\note{p. 51 de l'auteur.}
qui sont induits par les morphismes de normalisation entre multiplicités
associées aux graphes \uncertain{marqués}. Bien entendu, \add{pour \ill{} \uncertain{vraiment}}
\struck{\ill{}} \uncertain{bel et bien} les $\Pi_1 M$ des morphismes de normalisation, et
leurs propriétés de commutation aux \uncertain{bornes}, il faut utiliser la
\struck{connaissance} commutation de (148) à
\[
  \mathfrak{S}_{I,I'} \simeq \mathfrak{S}_{I'} \times \mathfrak{S}_{I\setminus I'} \;\ldots
\]

\emph{Programme}. Déterminer (à équivalence \uncertain{canonique} près) tous les groupoïdes
$\Pi_1 M_{g,I}$, avec leur dépendance fonctorielle pour $g$ fixé, $I$ variable
(des mono.\ $I' \hookrightarrow I$) via les morphismes de normalisation (148).
\marginal{par « générateurs et relations » en termes purement \uncertain{géométriques} …}

De plus, pour \struck{\ill{}} un MD-graphe $G$, déterminer l'extension
$S\Pi_1 M_G$ de
\[
  \Pi_1 M_G \simeq \prod_{\alpha\in S(G)} \Pi_1 M_{\alpha, \hat{I}_\alpha} \quad\text{par}\quad \mathbb{Z}^{\tilde{A}(G)},
\]
\note{il écrit ici $M_{\alpha, \hat{I}_\alpha}$, sans le $g$ des pages précédentes.}
et pour un morphisme de généralisation \struck{\ill{}} $G \to G'$ de MD-graphes
(p.~ex.\ un iso !) \struck{définir} déterminer les morphismes de groupoïdes de
généralisation correspondants $S\Pi_1 M_G \to S\Pi_1 M_{G'}$.

\page{104}
\note{toute la page est barrée d'un trait oblique ; elle est transcrite en entier. Le n° 12 est repris à la page suivante.}
\struck{12) \emph{Relation entre les morphismes de généralisation et morphismes de normalisation}.}
\note{dans le titre, « les » est biffé.}

\struck{Avant d'aborder une réalisation calculatoire de ce programme, j'ai envie
d'examiner les relations entre les deux types de morphismes, sur les
$\Pi_1 M_G$, $S\Pi_1 M_G$} --- \struck{\ill{}} \add{en dehors de}
\struck{la constatation $\pm$ triviale que, pour des isomorphismes de graphes,
les deux sont « compatibles » avec les projections}
\struck{$S\Pi_1 M_G \to \Pi_1 M_G$.} \struck{Pour le dire autrement, j'ai envie,
pour une inclusion \emph{stricte} $I' \subset I$, d'examiner « l'effet à
l'infini » des morphismes de groupoïdes (148).}

\struck{Soit donc $G_{g,I}$ le} \add{MD-}\struck{graphe réduit à un seul sommet de
genre $g$, et avec $I$ comme ensemble de pts marqués. Soit $G$ un graphe du type
numérique ($g$, $\nu = \operatorname{card} I$), et considérons un morphisme de
généralisation}
\[
  \text{(149)}\qquad \text{\struck{$G \to G_{g,I}$,}}
  \qquad \text{\struck{i.e.\ un iso.\ $I(G) \simeq I$}}
\]
\struck{d'où un diagramme de multiplicités}

\begin{tikzcd}[column sep=small, row sep=small]
  M^{*}_{G, G_{g,I}} \arrow[r, hook] \arrow[dr] & {(M_{G, G_{g,I}}, \Gamma^{!}_{G}) \simeq (M_G, \Gamma^{!}_G)} \\
  & M_{g,I}
\end{tikzcd}

\note{ce diagramme porte le numéro (150) ; à côté du numéro, un $M^{*}_G$ biffé. Le $\simeq (M_G, \Gamma^{!}_G)$ est écrit sous le terme de droite.}

\page{105}
\note{p. 52 de l'auteur. Le numéro du paragraphe se lit « 11) » ou « 12) », un chiffre récrit sur l'autre.}
\section{11) Digression : déploiement d'un diviseur à croisements normaux}

Pour fixer les idées, je me place dans le contexte des multiplicités analytiques
lisses --- mais on pourrait aussi bien, en prenant des multiplicités analytiques
\emph{relatives} --- ou en remplaçant les esp.\ analytiques par des esp.\ topologiques
ad-hoc --- ou en passant au contexte \uncertain{schématique} …

Soit $X$ une multiplicité analytique lisse, et $D$ un diviseur sur $X$ à
croisements normaux (hyp.\ \ill{} \ill{} sur $(X, D)$, par définition).
\struck{Soit $G$ un} Pour tout pt $x$ de $D$, l'ens.\ des « branches » de $D$
passant par $x$ est un ensemble fini $I(x)$.
\marginal{il conviendrait de prendre aussi, pour $x \in X$, et $x \notin D$, $I(x) = \emptyset$ (lecture incertaine)}
Par définition, les \add{germes de} sous-multiplicités \uncertain{attachées} à $D$ en $x$
sont les intersections d'un ens.\ fini \struck{de} parmi les $D_{x,i}$ ($i \in I(x)$),
\add{$\Delta_x$ ($= \Delta_x(X)$)}\note{lecture incertaine de cet ajout.}
les germes \struck{\ill{}} forment un ensemble ordonné \struck{Si} \add{le complémentaire}
anti-isomorphe à $\mathfrak{P}(I(x))$, ($J \subset I(x)$) étant \uncertain{associé} \struck{\ill{}}
à $D_{x,J}$,
\[
  \text{(149)}\qquad \operatorname{Codim}(D_{x,J}, X_x) = \operatorname{card} J
\]
\note{le numéro (149) est le même que celui de la page barrée précédente.}

\page{106}
Soit $J$ un ensemble ordonné, on s'intéresse aux plongements d'ensembles ordonnés
\[
  \text{(150)}\qquad J \overset{\varphi}{\hookrightarrow} \Delta_x \simeq \mathfrak{P}(I(x))^{\circ}
\]
Un tel plongement donne une « fonction codimension » sur $J$, un morphisme,
\add{soit $d$,} \struck{et} \struck{disons que $\varphi$ est un plongement d'un}
\uncertain{en supposant} donnée déjà $d$ \uncertain{dimension} sur $J$ \struck{ordonné, signifie que}
satisfaisant
\[
  \text{(151)}\qquad \text{si } x < y \text{ alors } c(x) < c(y),
\]
\note{le $c$ de gauche est récrit sur une autre lettre ($d$ ?).}
et on s'intéresse aux $\varphi$ qui sont des plongements ordonnés, compatibles aux
fonctions codim.\ sur $J$ et sur $\Delta_x$. On suppose de plus que $J$ a un plus
petit élément \struck{(qui correspond donc à une codim.\ minimale, soit $d(J)$) --- quitte à le rajouter, s'il n'y est déjà !}
\note{le passage entre parenthèses est encadré d'un trait vertical à gauche et barré de trois traits obliques.}

Considérons, \add{pour $p \in \mathbb{N}^{*}$,} la multiplicité analytique formée
des couples $(x, V)$, où $x$ est un pt de $D$ et $V$ un germe de sous-multiplicité
de $X$ attachée à $D$, de codim.\ $p$. On \struck{voit} écrit \uncertain{aisément} que ce sont là
les points d'une multiplicité analytique \uncertain{qui prolonge} $D_p$, et on a un morphisme

\page{107}
\note{p. 53 de l'auteur. L'indice de $D_d$, ici et à la fin de la page précédente, est récrit : un $d$ sur un $p$ ; on lit $d$.}
canonique
\[
  D_d \longrightarrow X
\]
qui se factorise d'ailleurs par $D$, et qui est un morphisme d'immersion locale,
de codim.\ $d$. Pour $D_d$, on a un revêtement étale \ill{} \ill{}, noté
$\widetilde{D}_d$, \struck{formé des couples} \add{dont la fibre en} un pt $(x, V)$ est
formée des \ill{} branches $D_{x,i}$ de $D$ en $x$ qui contiennent $V$. Donc les
pts de $\widetilde{D}_d$ correspondent aux triples $(x, V, W)$, où $x$ est un pt
de $D$, et $V \subset W$ des germes de ss-multiplicités de $X$ en $x$ attachées
à $D$, de codim.\ respectives $d$ et $1$. On trouve donc un morphisme canonique
\[
  \widetilde{D}_d \longrightarrow D_1, \qquad (x, V, W) \longmapsto W .
\]
\marginal{On a évidemment aussi : $D_0 = X$, (152) et $D_0 \to X$ l'identité.}

Plus généralement, si $J$ est un ens.\ \add{ordonné} muni d'une fonction codim.,
comme ci-dessus, \add{on en désigne} par $D_J$ la multiplicité \struck{formée}
\add{dont les} \uncertain{points}, de codim.\ $d = d(J)$, \add{et} sont les couples
d'un $x \in D$ \note{« $X$ » corrigé en « $D$ ».} et d'un plong.\ d'ens.\ ordonnés,
compatible avec fonctions codim.
\[
  \text{(152)}\qquad J \overset{\varphi}{\hookrightarrow} \Delta_x \simeq \mathfrak{P}(I(x))^{\circ}
\]
\marginal{on suppose la fonction codim.\ strictement \uncertain{croissante}, cf.\ (151)}
\note{le numéro (152) figure deux fois sur la page, dans la marge et ici.}

\page{108}
\struck{On a \struck{déjà} supposé que $J$ a un plus petit élément $e_J$, dont $d(J) = d$ sa codim.}
\note{ces deux lignes sont encadrées et barrées.}

Alors la multiplicité $D_J$ est lisse, \struck{\ill{}} et le morphisme canonique
\[
  \text{(153)}\qquad D_J \longrightarrow X
\]
est une imm.\ locale de codim.\ $d$. Elle se factorise par $D$, sauf dans le
cas où $d(J) = 0$ \add{i.e.\ $J$} \struck{est} réduit à un seul élément, de codim.\ nulle,
auquel cas $D_J = D_0 = X$. Si $J$ est formé d'un seul élément, de codim.\ $d$,
on retrouve $D_d$ ; s'il est formé de deux éléments, \struck{le plus} $e_J$ de
codim.\ $d$ et $e'$ de codim.\ $1$, on retrouve $\widetilde{D}_d$.

Si $J$ est de \ill{} \uncertain{totale} $d = d(J)$, \uncertain{on peut} \struck{\ill{}} reconstruire
$D_J$ à partir de $D_d$ et de $\widetilde{D}_d$, via le \uncertain{revêtement} principal
(\uncertain{tordu} sous $\mathfrak{S}_d$) \uncertain{associé} à $\widetilde{D}_d$, soit
$D^{!}_d$, qui \add{au-dessus de} $D_d$ s'identifie aussi à $D_{I(d)}$, où $I(d)$
\uncertain{désigne} l'ens.\ ordonné

\begin{tikzcd}[column sep=small]
  d \arrow[r, no head] & d-1 \arrow[r, no head] & 2 \arrow[r, no head] & 1 \arrow[r, no head] & 0
\end{tikzcd}

(drapeau \uncertain{maximal} de codim.\ $d$). En effet, considérons,
\marginal{Mais l'expression de $D_{I(d)}$ en termes de $D_d$ \ill{} \ill{} $D^{!}_d$ \ill{} \ill{} $\mathfrak{S}_d$ \ill{} \ill{} le plus petit élément \ill{} les codim.\ maximales \ill{} (note écrite en oblique dans la marge gauche, en grande partie illisible)}
\note{le drapeau est dessiné comme une suite de points reliés, codimensions décroissantes de $d$ à $0$.}

\page{109}
\note{p. 54 de l'auteur.}
l'ensemble fini
\[
  \text{(154)}\qquad H(J) = \text{ensemble des plongements d'ens.\ ordonnés,}
\]
compatibles avec les f.\ codim., de $J$ dans \struck{$\mathfrak{P}([1,n]\cap\mathbb{N})$}
$\mathfrak{P}(I_d)^{\circ}$ (où $I_d \overset{\text{déf}}{=} [1,d] \cap \mathbb{N}$) ;
alors $\mathfrak{S}_d$ opère à gauche sur $H(J)$, et on a un iso.\ can.\ de
$D$-multiplicités
\[
  \text{(155)}\qquad D_J \simeq D^{!}_d \overset{(\mathfrak{S}_d)_{D_d}}{\wedge} H(J) .
\]
\note{au-dessus de $D^{!}_d$, une marque illisible.}

Soient \struck{\ill{}} $(J, \underline{c})$ et $(J', \underline{c}')$ deux ens.\ ordonnés finis munis de
fonctions codim., et
\[
  \text{(156)}\qquad J' \overset{f}{\longrightarrow} J
\]
un plongement d'ens.\ ordonnés, compatible aux fonctions codim., on trouve alors
un \struck{application} morphisme canonique
\[
  \text{(157)}\qquad D_J \longrightarrow D_{J'} ,
\]
\uncertain{qui se} \ill{} \struck{\ill{}} \struck{en termes de} quelles données supplémentaires
\add{(si besoin est)} sur les $(D_d, D^{!}_d)$ \uncertain{peut-on} reconstituer cette
dépendance fonctorielle \add{comme la} \struck{considérée} de $D_J$ p.r.\ à $J$, via (155).
\struck{d'abord le plus petit élément $e_{J'}$ de $e$} \ill{} dépendance fonctorielle en $J$
pour les \emph{iso}

\page{110}
est claire sur l'expression (155) ; soit $J' \subset J$. Considérons
\struck{le plus petit} \add{un} élément \add{$x$} \struck{$e_J$} de $J'$, $\neq e_J$,
\struck{élément} de codim.\ $\delta$ (\add{$< d$}), et considérons \add{d'abord} le composé
\[
  \text{(158)}\qquad D_J \longrightarrow D_{J'} \overset{f_x}{\longrightarrow} D_\delta
\]
\add{défini par \uncertain{tautologie} $\{x\} \subset J'$}
\note{sous la flèche, un passage biffé : « $D_{\{x\}}$ \ill{} iso. ».}
on va la mettre en relation avec l'homm.
\[
  \text{(159)}\qquad D_J \longrightarrow D_{d,\delta}
\]
suivant (où $D_{d,\delta} = D_{\Delta(d,\delta)}$), $\Delta(d,\delta)$ étant le drapeau
\struck{type} du type $d, \delta$, i.e.\ l'ens.\ ordonné \struck{avec} à fonction codim.

\begin{tikzcd}
  d \arrow[r, no head] & \delta
\end{tikzcd}

(car $\delta < d$ !) \struck{\ill{} $d = \delta$} \struck{(car $\delta < d$)}. \uncertain{Utilisant} \ill{},
on peut prendre (155) avec $J = \Delta(d,\delta)$, \struck{au lieu de} \ill{} l'isom.\ canonique
\[
  \text{(160)}\qquad D_{d,\delta} \simeq D^{!}_d \overset{(\mathfrak{S}_d)_{D_d}}{\wedge} \mathfrak{P}_\delta(I_d)
\]
D'autre part, on a (dans la notation (154) pour $H(J)$), pour l'élément $x \in J$
fixé ($\delta$ \uncertain{$= \operatorname{cod}(x)$}), une application
\[
  \text{(161)}\qquad H(J) \longrightarrow \mathfrak{P}_\delta(I_d), \qquad h \longmapsto h(x)
\]
\note{la lettre $x$ est récrite sur une autre, ici comme plus haut.}

\page{111}
\note{p. 55 de l'auteur.}
compatible avec les opérations de $\mathfrak{S}_d$, d'où un morphisme

\begin{tikzcd}
  D_J \arrow[r, "f_x"] \arrow[d, no head, "\wr"'] & D_{d,\delta} \arrow[d, no head, "\wr"] \\
  {\widetilde{D}_d \wedge_{\mathfrak{S}_d} H(J)} & {D^{!}_d \wedge_{\mathfrak{S}_d} \mathfrak{P}_\delta(I_d)}
\end{tikzcd}

\note{ce diagramme porte le numéro (162) ; les deux $\simeq$ sont écrits verticalement, sans flèche.}
\struck{Ceci posé,} Cette relation est décrite entièrement en termes du rev.\
principal $\widetilde{D}_d$ de $D_d$, de groupe $\mathfrak{S}_d$, \add{et de $x \in J$.}
On \add{a « un $J$- de $D_J$ »} \note{lecture incertaine de cet ajout.} d'autre part, par
« fonctorialité », la \struck{\ill{}} commutativité \uncertain{étendant} (158)
\[
  \text{\struck{(162)\quad $D_{d,\delta} \to D_{\delta'}$}}
\]
\struck{et on a}

\begin{tikzcd}
  D_J \arrow[r] \arrow[d] & D_{J'} \arrow[d, "f_x"] \\
  D_{d,\delta} \arrow[r] & D_\delta
\end{tikzcd}

\note{ce carré porte le numéro (163).}
où $D_J \to D_{d,\delta}$ et $D_{d,\delta} \to D_\delta$ sont définis par
fonctorialité par les inclusions $\{x\} \subset \{x, e\} \subset J$. On
\struck{a l'impression} \add{\ill{}} que $D_J \to D_{d,\delta}$ est décrit par la relation
(162), donc le composé (158) est connu quand on connaît les $D_{d,\delta} \to D_\delta$

\page{112}
--- et quand les composés (158) sont connus, on vérifie que les morphismes
$D_J \to D_{J'}$ sont connus …

On est donc intéressé plus particulièrement par les homs
\[
  \text{(164)}\qquad D_{d,\delta} \longrightarrow D_\delta \qquad (d > \delta)
\]
qui jouent le rôle des « morphismes but », alors que les $D_{d,\delta} \to D_d$
jouent le rôle de « morphismes source », et $D_{d,\delta}$ le rôle de
« \struck{\ill{}} multiplicité des flèches » (\uncertain{ici} morph.\ d'inclusion) --- entre
éléments de la sous-multiplicité \add{de $X$} associée à $D$.
On a des isom.\ can.
\[
  \text{(165)}\qquad D_{d,d'} \times_{D_{d'}} D_{d',d''} \simeq D_{d,d',d''}
  \qquad \text{si } d > d' > d''
\]
et on en déduit de cette description les morphismes canoniques

\begin{tikzcd}[column sep=small, row sep=tiny]
  & D_{d,d'} \arrow[r] \arrow[dr] & D_d \\
  D_{d,d',d''} \arrow[ur] \arrow[dr] & & D_{d'} \\
  & D_{d',d''} \arrow[ur] \arrow[r] & D_{d''}
\end{tikzcd}

\note{ce diagramme porte le numéro (166).}
mais non directement \struck{$D_{d'}$}
\[
  \text{(167)}\qquad D_{d,d',d''} \longrightarrow D_{d,d''},
\]
mais on sait que
\[
  \text{(168)}\qquad D_{d,d''} \longrightarrow D_d \times D_{d''}
\]

\page{113}
\note{p. 56 de l'auteur.}
est un \uncertain{mono}, et (167) est déterminé par la condition que son composé avec
(168) est donné par les morphismes
\[
  D_{d,d',d''} \rightrightarrows D_d,\ D_{d''}
\]
qui figurent dans (166).
\marginal{ceci marche dans le contexte des esp.\ analytiques (lecture incertaine)}
\note{dans les renvois de cette page, le dernier chiffre de (167) et (168) est récrit.}

Une « structure de déploiement d'intersection \emph{normale} » (d'abord dans le
contexte des \emph{espaces} analytiques) sera définie \add{d'abord} comme une
\struck{multiplicité} structure de « catégorie analytique » $\mathcal{D}$, i.e.\ d'une
catégorie dans la catégorie des espaces analytiques :
\marginal{on pourrait dire « structure de déploiement analytique » (lecture incertaine)}
\[
  \text{(169)}\quad
  \begin{cases}
    D_{*} & \text{espace analytique des objets} \\
    D_{**} & \text{\phantom{espace analytique des} flèches}
  \end{cases}
\]

\begin{tikzcd}
  & D_{**} \arrow[dl, "s"'] \arrow[dr, "b"] & \\
  D_{*} & & D_{*}
\end{tikzcd}

morphismes « source » et « but » ;
\[
  D_{*} \longrightarrow D_{**} \qquad \text{morphisme « identité »} ;
\]
\[
  (D_{**}, b) \times_{D_{*}} (D_{**}, s) \longrightarrow D_{**}, \qquad (u, v) \longmapsto v \circ u
\]
(morphisme « composition des flèches »)
\note{les quatre données sont réunies par une accolade sous le numéro (169).}

satisfaisant aux conditions \struck{\ill{}} \uncertain{habituelles} (associativité, unités …)
exprimant que

\page{114}
pour tout $S$, les « pts » de $\mathcal{D}$ à valeurs dans $S$ (formés de
$D_{*}(S)$, $D_{**}(S)$ etc.)\ est une catégorie. On suppose de plus donnée une
fonction « codimension » sur $D_{*}$
\[
  \text{(170)}\qquad D_{*} \overset{c}{\longrightarrow} \mathbb{N}
\]
\marginal{c'est inutile, elle résulte des autres structures, \uncertain{cf.}\ plus loin … (lecture incertaine)}
ce qui revient à une décomposition
\[
  \text{(171)}\qquad D_{*} = \coprod_{d\in\mathbb{N}} D_d
\]
d'où une décomposition correspondante
\[
  \text{(172)}\qquad D_{**} = \coprod_{(d,d')\in\mathbb{N}\times\mathbb{N}} D_{d,d'}
\]
où
\[
  \text{(173)}\qquad D_{d,d'} = \text{image inverse de } D_d \times D_{d'}
  \text{ par } D_{**} \xrightarrow{(s,b)} D_{*} \times D_{*} .
\]
On fait les hypothèses suivantes
\note{les hypothèses a) à c) sont réunies par une accolade sous le numéro (174).}
\begin{enumerate}
  \item[a)] $D_{d,d'} = \emptyset$ si $d' > d$ (le but est de codim.\ $\leq$ la codim.\ de la source) ;
  \item[b)] le morphisme « identité » $D_{*} \to D_{**}$ \struck{induit} (qui se factorise en
  $\coprod_{d\in\mathbb{N}} D_d \to \coprod_{d\in\mathbb{N}} D_{d,d}$) induit un iso
  \[
    D_d \simeq D_{d,d}
  \]
  (une flèche dont source et but sont de $=$ codim., est une identité) ;
  \marginal{ce sera conséquence de c) ci-dessous (\ill{} \ill{})}
  \item[c)] Pour $d \in \mathbb{N}$ fixé, le morphisme « source » \struck{$D_{d*} \to D_{*}$}
\end{enumerate}
\note{la phrase se poursuit à la page suivante.}

\page{115}
\note{p. 57 de l'auteur. Suite de l'hypothèse c) de (174), le numéro (174) étant répété en marge.}
\[
  \Bigl(\; D_{d,*} \overset{s_d}{\longrightarrow} D_d, \qquad
  D_{d,*} \overset{\text{déf}}{=} \coprod_{d'\in\mathbb{N}} D_{d,d'} = s^{-1}(D_d) \;\Bigr)
\]
qui fait de $D_{d,*}$ un esp.\ analytique sur $D_d$, \struck{\ill{}} relativement ordonné
(par la relation de factorisabilité des flèches de source donnée
\drawing{petit schéma : un point d'où partent deux flèches, la seconde se factorisant par la première au moyen d'une flèche en pointillé ; annoté « prédicat \uncertain{par} $d'$ »})
est loc.\ isomorphe \add{avec sa relation d'ordre et la fonction « codim. »} à
$\mathfrak{P}(I_d)_{D_d}$ --- le faisceau constant d'ensembles ordonnés
\uncertain{simpliciaux} $\mathfrak{P}(I_d)^{\circ}$ ---

Cela redonne donc un rev.\ $D^{!}_d$ principal de groupe $\mathfrak{S}_d$, permettant
de reconstituer les $D_{d,d'}$ --- mais pour reconstituer $\mathcal{D}$, il faut de
plus se donner \struck{\ill{}} --- \uncertain{mieux}, les $D_{d,d'} \to D_{d'}$, et (c'est plus
sûr !)\ aussi
\[
  D_{d,d'} \times_{D_{d'}} D_{d',d''} \longrightarrow D_{d,d''}
\]
(les flèches composition), car il n'est pas clair, avec les seules hyp.\ faites
jusqu'à présent, que le morphisme
\[
  D_{d,d''} \longrightarrow D_d \times D_{d''}
\]
soit mono., i.e.\ que « les flèches sont déterminées par leur source et par leur
but ». Mais pourtant si, c'est une conséquence de c). Donc

\page{116}
la structure envisagée peut se décrire par
\begin{enumerate}
  \item[a)] famille d'espaces analytiques $(D_d)_{d\in\mathbb{N}^{*}}$
  \item[b)] $\forall\, d \in \mathbb{N}$, un revêtement principal $D^{!}_d$ de $D_d$, de
  groupe $\mathfrak{S}_d$ --- d'où des $D_{d,d'}$ ($0 \leq d' \leq d$).
  \item[c)] des morphismes
  \[
    D_{d,d'} \overset{b}{\longrightarrow} D_{d'} \qquad (\text{pour } 0 \leq d' < d),
  \]
  et en plus des diagrammes cartésiens
\end{enumerate}
\marginal{--- on a $D_{d,d} \simeq D_{d,0} \simeq D_d$ ; $D_{d,1}$ est le rev.\ de $D_d$ à $d$ feuillets associé à $D^{!}_d$ (lecture incertaine)}
\struck{la condition que les morphismes can.\ (pour $d > d' > d''$)}
\[
  \text{\struck{$D_{d,d',d''} \overset{\text{déf}}{=} D_{d,d'} \times_{D_{d'}} D_{d',d''} \to D_d \times D_{d''}$}}
\]
\note{ce passage, ainsi que deux débuts de d) (« d) $D_{d,d'} \times_{D_{d'}} D^{!}_{d'} \simeq$ … », « d) $D^{!}_{d,d'}$ … »), est biffé de traits croisés.}

d)

\begin{tikzcd}
  D^{!}_{d,d'} \arrow[r] \arrow[d, "\mathrm{can}"'] & D^{!}_{d'} \arrow[d, "\mathrm{can}"] \\
  D_{d,d'} \arrow[r, "b"'] & D_{d'}
\end{tikzcd}

\note{ce carré porte le numéro (175) et la mention « cart. » au centre.}
où $D^{!}_{d,d'}$ est le revêtement $\mathfrak{S}_{d'}$-principal de
\[
  D_{d,d'} \simeq D^{!}_d \overset{(\mathfrak{S}_d)_{D_d}}{\wedge} \mathfrak{P}_{d'}(I_d),
\]
associé au $d'$-revêtement « tautologique » de $D_{d,d'}$) --- ce qui
\uncertain{se déduit de ce} donné, puisque $D_{d,d',d''}$ rev.\ de $D_{d,d'}$ par
\[
  D_{d,d',d''} \overset{\text{déf}}{=} D^{!}_{d,d'} \overset{(\mathfrak{S}_{d'})_{D_{d,d'}}}{\wedge} \mathfrak{P}_{d''}(I_{d'}),
\]
un \struck{\ill{}} diagramme cartésien

\begin{tikzcd}
  D_{d,d',d''} \arrow[r] \arrow[d] & D_{d',d''} \arrow[d] \\
  D_{d,d'} \arrow[r] & D_{d'}
\end{tikzcd}

\note{ce carré porte le numéro (176) et la mention « cart. » au centre.}
\marginal{NB $D_{d,d'}$ et $D_{d,d',d''}$ se déduisent \uncertain{directement} de $D^{!}_d$}

\page{117}
\note{p. 58 de l'auteur.}
et il faut demander de plus qu'on ait commutativité dans les \struck{morph.}\
diagrammes

\begin{tikzcd}
  D_{d,d',d''} \arrow[r] \arrow[d, "\text{tautol.}"'] & D_{d',d''} \arrow[d, "\text{« but »}"] \\
  D_{d,d''} \arrow[r, "\text{« but »}"'] & D_{d''}
\end{tikzcd}

\note{ce carré porte le numéro (177).}

\emph{Remarque}. L'axiomatique précédente a un sens dans \add{un} site
(\struck{\ill{}} le fait de \ill{} un \ill{} dans un site donné un sens à
\ill{} des \ill{} de revêtement principal …), où les fonctions loc.\ constantes
\uncertain{à valeurs} \ill{} définissent une décomposition en somme disjointe d'un objet.
Il est inutile de se \emph{donner} la fonction codimension, elle est contenue dans
la condition c), qui implique a) et b).

\struck{Pour} \uncertain{Pour} faire le lien avec l'intuition des « croisements normaux », il
faut faire des hypothèses ad-hoc sur la nature des morphismes
\[
  D_{d,d'} \longrightarrow D_{d'}
\]
(NB $D_{d,d'} \to D_d$ est un rev.\ fini étale …), p.~ex.

\page{118}
que c'est une immersion \emph{locale} d'espaces \struck{\ill{}} analytiques lisses (ou
lisses sur un esp.\ analytique donné, dans une situation relative …) ;
\struck{de codim.\ $d - d'$} \struck{\ill{} \ill{} fait l'hyp.\ de}
\struck{Notons que pour que} \struck{\ill{},} \struck{$D_{d,d'} \to D_{d'}$ soit une imm.\ locale, il suffit}
\struck{que $D_{d,0} \to D_0$ et $D_{d',0} \to D_0$} \struck{($D_{d,0} \simeq D_d$, $D_{d',0} \simeq D_{d'}$) le soient, comme on voit en considérant}
\note{tout ce passage, encadré, est barré de traits obliques. Il se termine par un diagramme, barré lui aussi, où $D_{d,d',0} \to D_{d',0} \to D_0$ et $D_{d,d',0} \to D_{d,0} \to D_0$ se placent au-dessus d'un carré cartésien de base $D_{d,d'} \to D_{d'}$.}

Dans le cas lisse ou relatif lisse, l'hyp.\ d'imm.\ locale équivaut d'ailleurs à
celle de mono.\ local, qui garde un sens dans tout site.

Notons que pour tout $d$ on a

\begin{tikzcd}
  D_{d,0} \arrow[d, no head, "\wr"'] \arrow[dr] & \\
  D_d & D_0 \overset{\text{déf}}{=} X
\end{tikzcd}

\note{ce diagramme porte le numéro (178).}
donc les $D_d$ s'envoient tous canoniquement dans $X = D_0$, et on aura
commutativité des

\begin{tikzcd}
  D_{d,d'} \arrow[r, "b"] \arrow[d, "s"'] & D_{d'} \arrow[d] \\
  D_d \arrow[r] & X
\end{tikzcd}

\note{ce carré porte le numéro (179).}
\marginal{$\mathcal{D}$ est une catégorie dans la catégorie (des objets) au-dessus de $D_0$}

\page{119}
\note{p. 59 de l'auteur.}
où $s$ est un rev.\ étale donc un iso.\ local. Donc on voit que si
$D_d \to D_0$ et $D_{d'} \to D_0$ sont des immersions locales (resp.\ des mono.\
locaux), il en est de même de $D_{d,d'} \to D_{d'}$. De plus, dans le cas lisse,
on voit que la codim.\ de $D_{d,d'} \to D_{d'}$ est égale à la différence des
codim.\ de $D_d \to D_0$ et $D_{d'} \to D_0$. On supposera, bien sûr (condition
\uncertain{essentielle} ! (?))
\[
  \text{(180)}\qquad D_d \longrightarrow D_0 \text{ est de codim.\ } d
\]
\struck{NB} et plus précisément encore ceci. Soit $x \in D_d$, d'où $I(x)$ de
cardinal $d$ (tel qu'on ait
\[
  \text{(181)}\qquad (D^{!}_d)_x \simeq \mathfrak{P}(I_x) \;)
\]
\note{le symbole devant $(I_x)$ est celui qu'il emploie pour $\mathfrak{P}$ ; lecture incertaine ici.}
\struck{\ill{} considérons les} \ill{}
\[
  \text{(18\ill{})}\qquad I_x = \text{\uncertain{ensemble} des sections en $x$ du rev.\ $D_{d,1}$ au-dessus de $D_d$}
\]
\[
  \simeq (D_{d,1})_x \;\text{---},
\]
Considérons alors les composés ($i \in I$)
\[
  D_d \overset{s_i}{\longrightarrow} D_{d,1} \overset{\mathrm{can}}{\longrightarrow} D_1 \longrightarrow X
\]
qui sont tous égaux au morph.\ can.\ $D_d \overset{\mathrm{id}}{\longrightarrow} X$,
\note{dans ces deux formules, une lettre récrite ou biffée avant $X$.}
d'autre part les $x_i$ que \uncertain{définit} $x$ dans $D_1$, pr.\ les $i$.

\page{120}
définissent des germes de \uncertain{sous-espaces} \add{$D^{i}_x$} de $X$ en l'image
de $x$ dans $X$, qui contiennent le germe défini par $(D_d, x) \to (X, a)$
(ceci dit, on veut que ce dernier germe soit l'\emph{intersection} des germes
$D^{i}_x$, \add{\uncertain{et}} \ill{} que cette intersection soit \emph{transversale}).

Cette axiomatique inclut des cas tels que
\drawing{courbe plane fermée qui se recoupe elle-même en une boucle}
\[
  X = D_0 \simeq E^2, \qquad D_1 = \text{\struck{\ill{}}\ \emph{lisse} conn.}, \qquad D_2 = \emptyset \;\ldots
\]
\marginal{(normalisée de la courbe \uncertain{tracée})}
qui ne peut être décrit par la donnée d'un diviseur sur $X$,
\struck{\ill{} \ill{}} \struck{le diviseur, donc}

L'opération de $\mathfrak{S}_d$ sur $\mathbb{Z}^d$ permet de construire sur $D_d$
un système local
\[
  \text{(183)}\qquad T_d = D^{!}_d \overset{(\mathfrak{S}_d)_{D_d}}{\wedge} (\mathbb{Z}^d)_{D_d} ,
\]
qui donne naissance à un groupe analytique (ou algébrique) relatif sur $D_d$ :
\[
  \text{(184)}\qquad G_d = T_d \otimes_{\mathbb{Z}} \mathbb{G}_{m, D_d}
  \simeq D^{!}_d \overset{(\mathfrak{S}_d)_{D_d}}{\wedge} \mathbb{G}_m^d
\]
\struck{\ill{}} et $T_d$ s'interprète en termes de $G_d$ comme le syst.\ local des
$\pi_1$ ou $H_1$ des fibres
\note{la phrase se poursuit au-delà de la fin du lot.}

\end{document}
